\section{실험 결과}
\label{sec:results}

\begin{table}[t]
\centering
\bitablecaption{중심 결과와 직접 지지되는 주장.}{Primary results and directly supported claims.}
\label{tab:core-results}
\footnotesize
\begin{tabularx}{\columnwidth}{L{0.27\columnwidth}L{0.34\columnwidth}Y}
\toprule
Result & Supported claim & Not supported \\
\midrule
40 synthetic pairs: prior-requirement checker found 40/40 modified cases; current-only found 20/40; both found 0/40 reference cases & Retaining prior requirements finds both cross-event error types; all four test rules were exercised & Accuracy on unseen natural errors or vehicle safety \\
Python--UPPAAL mismatch 0/7 & Both implementations give the same decisions and reasons on tested cases & Independent proof of the rules or vehicle safety \\
27 selected scenes: textual consistency disputed in 23/27; $\kappa=0.027$; contradiction 0/27 after discussion & Ambiguity shows the need for reviewer discussion; 0/27 is the selected-group record & Natural contradiction rate or detection accuracy \\
\bottomrule
\end{tabularx}
\end{table}

\begin{table}[t]
\centering
\bitablecaption{보조 결과와 직접 지지되는 주장.}{Auxiliary results and directly supported claims.}
\label{tab:aux-results}
\footnotesize
\begin{tabularx}{\columnwidth}{L{0.28\columnwidth}L{0.34\columnwidth}Y}
\toprule
Result & Supported claim & Not supported \\
\midrule
\texttt{scene-0001}: reference case passes; four modified cases fail & Checker distinguishes four deliberately added errors & Full-dataset performance or natural error rate \\
$134=125+4+5$ at $D=3$ s; two decisions reverse & Deadline and repeated-event rules change results & Safety or the best repeat rule \\
130 targets/27 settings: 72 unchanged, 58 varied & Speed-based evidence depends on its settings & Physical validity of every setting \\
Correct boundary handling passes; faulty carryover fails & Unsupported transfer to the next scene is blocked & Long-horizon or object-continuity accuracy \\
\bottomrule
\end{tabularx}
\end{table}

표~\ref{tab:core-results}와 표~\ref{tab:aux-results}는 각 수치가 어떤 결론을 뒷받침하고 어떤 결론까지는 뒷받침하지 않는지를 한 줄씩 보여 준다. 먼저 합성 비교로 중심 주장을 확인하고 Python과 UPPAAL의 결과가 같은지 살핀다. 이어 실제 주행 자료에서 여러 사람이 주석을 얼마나 비슷하게 해석하는지 확인하고, 반응 시간 설정과 장면 경계 처리에 따라 결과가 달라지는 범위를 설명한다.

\subsection{이전 요구 기억 여부를 비교한 합성 시험}

표~\ref{tab:sequential-mutation-results}에서 이전 요구를 기억하는 검사는 바꾼 사례 40건 모두에서 미리 지정한 오류를 판정했고, 현재 사건만 보는 검사는 20건에서만 오류를 판정했다. 두 검사 모두 정상 기준 사례 40건에는 오류를 보고하지 않았다. 차이는 이전 사건 정보가 꼭 필요한 두 경우에서만 나타났다. 하나는 정지ㆍ양보 요구가 끝나기 전에 진행한 경우이고, 다른 하나는 이미 끝난 요구를 검사 목록에서 지우지 않은 경우이다. 이 두 종류는 현재 사건만 보면 0/10, 이전 요구를 기억하면 10/10이었고, 현재 사건 정보만으로 판단할 수 있는 나머지 두 종류는 결과가 같았다.

\begin{table}[H]
\centering
\bitablecaption{합성 기준--변형 쌍에서 이전 요구 기억 여부에 따른 판정 결과.}{Results for synthetic reference--modified pairs with and without prior-requirement memory.}
\label{tab:sequential-mutation-results}
\scriptsize
\begin{tabularx}{\columnwidth}{L{0.43\columnwidth}rr}
\toprule
Consistency error & Current only & Prior retained \\
\midrule
Proceed while stop/yield remains & 0/10 & 10/10 \\
End a requirement too early & 10/10 & 10/10 \\
Violate the required action order & 10/10 & 10/10 \\
Keep a requirement after its end & 0/10 & 10/10 \\
\midrule
All modified cases & 20/40 & 40/40 \\
All reference cases & 0/40 & 0/40 \\
\bottomrule
\end{tabularx}
\end{table}

이 동일 입력 비교는 정지ㆍ양보 요구가 끝나기 전 진행한 오류와 끝난 요구를 목록에 남긴 오류를 찾으려면 이전 요구를 기억해야 함을 직접 보여 준다. 바꾼 사례를 만들 때 기대한 오류가 실제로 나오는지 미리 확인했으므로, 40/40은 정해 둔 네 규칙이 시험 사례에서 모두 실행됐다는 뜻이다. 처음 보는 자연 오류를 얼마나 정확히 찾는지는 별도로 평가해야 한다.

\subsection{UPPAAL 구현 일치}

다음으로 같은 검사 규칙을 Python과 UPPAAL에서 실행했을 때 결과가 같은지 확인했다. 실제 UPPAAL 검사 프로그램을 실행한 일곱 표준 시험 사례에서 최종 판정, 오류 종류, 첫 위반 사건과 {\unknown} 사유가 모두 일치하여 불일치는 0/7이었다. 정지 요구가 끝나기 전에 진행한 사례에서는 끝나지 않은 정지 요구가 목록에 남은 상태에서 진행 사건을 읽자, 그 사건을 첫 위반으로 기록했다. 따라서 두 구현은 시험한 사례에서 같은 판정과 위반 기록을 재현했다. 다만 둘 다 같은 규칙을 구현하므로, 이 결과가 규칙 자체의 옳음이나 가능한 모든 주행 환경의 안전을 증명하지는 않는다.

\subsection{자연 표본 검토}

\begin{table}[H]
\centering
\bitablecaption{27개 선택 장면에 대한 합의 전 검토자 일치도.}{Reviewer agreement before discussion on 27 selected scenes.}
\label{tab:reviewer-agreement}
\footnotesize
\begin{tabular}{lrr}
\toprule
Review item & Pair agreement & Fleiss $\kappa$ \\
\midrule
Temporal requirement & 0.698 & 0.203 \\
Required action sequence & 0.340 & 0.291 \\
Textual consistency & 0.407 & 0.027 \\
\bottomrule
\end{tabular}
\end{table}

표~\ref{tab:reviewer-agreement}처럼 독립 판정에서 텍스트 일관성의 검토자 쌍 일치 비율은 0.407이었다. Fleiss $\kappa$는 우연히 같은 답을 고를 가능성을 뺀 일치 정도로, 1에 가까울수록 의견이 잘 맞고 0에 가까울수록 거의 맞지 않는다. 이 값은 0.027이었으며, 텍스트 일관성 항목은 23/27장면에서 의견이 갈렸다. 세 검토 항목 중 하나 이상에서는 27/27장면 모두 의견 차이가 있었다. 근거를 함께 검토한 뒤에는 텍스트 일관 27/27, 모순 0/27로 합의했다. 이는 자동 주석을 사람마다 다르게 읽을 수 있어 합의 절차가 필요함을 보여 주며, 0/27은 선택한 장면에 대한 검토자 합의 기록이다.

\begin{table}[H]
\centering
\bitablecaption{합의된 문장 판정과 기록 궤적 판정의 장면별 교차표.}{Scene-level comparison of agreed text decisions and recorded-trajectory decisions.}
\label{tab:logic-trajectory-cross}
\footnotesize
\begin{tabular}{lrrr}
\toprule
Consensus text & Aligned & Not aligned & Unknown \\
\midrule
Consistent & 2 & 7 & 18 \\
Contradictory & 0 & 0 & 0 \\
\bottomrule
\end{tabular}
\end{table}

표~\ref{tab:logic-trajectory-cross}에서 문장들은 모두 서로 모순이 없다고 합의됐지만, 문장이 요구한 행동과 기록된 이동 경로의 관계는 2장면 일치, 7장면 불일치, 18장면 판단에 필요한 정보 부족({\unknown})이었다. 이는 문장끼리 모순이 없는지와 문장이 기록된 행동에 맞는지가 서로 다른 문제임을 보여 준다. 불일치 7건도 상대 객체가 요구 종료 조건을 만족했는지와 차량의 실제 움직임을 설명할 정보가 없으므로, 교통법규나 물리적 안전을 위반했다고 판단할 수 없다.

\subsection{보조 검사기의 의도적 변경 시험과 반응 설정 결과}

앞선 결과가 문장 사이의 모순에 답했다면, 다음 결과는 반응 시간과 자료 출처 처리 방식이 검토 결과를 어떻게 바꾸는지 보여 준다. \texttt{scene-0001}의 첫 감속 요구는 2.5초, 자차 속도에서 찾은 반응 시작은 4.6597초여서 약 2.1597초가 걸렸다. 표~\ref{tab:scene1-results}의 변경하지 않은 기준 모델은 제한 시간 $D=3$초에서 모든 검사 조건을 통과했다. 제한 시간을 2초로 줄이거나, 반복 사건에서 시계를 다시 시작하거나, 반응을 제거하거나, 낮은 품질 사건이 기존 정보를 덮어쓰게 했을 때는 각각 관련 검사 조건이 실패했다. 표에서 P는 통과(pass), F는 실패(fail)를 뜻하며, 여섯 열은 요구 상태, 반응 확인, 제한 시간, 반복 시계, 정보 덮어쓰기와 멈춤 없는 실행을 차례로 나타낸다.

\begin{table}[H]
\centering
\bitablecaption{\texttt{scene-0001}의 기준 모델과 한 항목씩 바꾼 모델의 검사 결과.}{Check results for the reference scene-0001 model and four one-item modifications.}
\label{tab:scene1-results}
\scriptsize
\setlength{\tabcolsep}{2.5pt}
\begin{tabular}{lcccccc}
\toprule
Case & Req. & Resp. & Time & Repeat & Source & Run \\
\midrule
reference & P & P & P & P & P & P \\
2 s limit & P & P & F & P & P & P \\
restart time & P & P & P & F & P & P \\
no response & P & F & F & P & P & P \\
rejected overwrite & P & P & P & P & F & P \\
\bottomrule
\end{tabular}
\end{table}

바꾼 내용과 관련된 검사 조건만 실패했으므로, 기준 모델이 단순히 검사를 꺼 두어서 통과한 것은 아니다. 즉 검사기는 의도적으로 넣은 네 오류를 구분했다. 이어 전체 방향성 사건에서 설정을 바꿀 때 판정이 얼마나 달라지는지 확인했다.

반응 시간 보조 분석에서 행동 방향이 분명한 134건은 이후 속도 변화를 찾아야 하는 130건과, 사건 시점에 이미 정지ㆍ양보 요구를 만족한 4건으로 구성된다($134=130+4$). 기본 제한 시간 3초와 첫 판단 시각부터 계속 시간을 재는 방식에서 130건 중 125건은 속도 반응을 찾았고 5건은 검토 후보였다. 이미 만족한 4건을 더하면 $134=125+4+5$이다. 처음에는 후보가 8건이었지만, 사건 시점 속도가 $v\leq0.5$~m/s인 정지ㆍ양보 요구는 이미 충족한 것으로 보는 탐색 규칙을 나중에 적용해 세 건을 다시 분류했다. 따라서 129건은 이 검사 설정을 통과한 항목이고, 다섯 항목은 추가 검토가 필요하다.

\begin{table}[H]
\centering
\bitablecaption{반복 사건의 시간을 재는 두 방식에서 다섯 검토 후보가 통과한 제한 시간.}{Passing deadlines for five review candidates under two ways of timing repeated events.}
\label{tab:candidate-summary}
\scriptsize
\begin{tabularx}{\columnwidth}{L{0.11\columnwidth}L{0.25\columnwidth}L{0.25\columnwidth}Y}
\toprule
Scene & Required action & Passing $D$ (s): first-event/restart-on-repeat & Observation \\
\midrule
0019 & Yield/decelerate & none / none & Response found in 4/27 settings \\
0028 & Red-light stop & 4,5 / 3,4,5 & Reversal at $D=3$ s \\
0042 & Lead-car deceleration & 4,5 / 4,5 & Changes at 3--4 s \\
0065 & Bus deceleration & 4,5 / 4,5 & Other-object state unavailable \\
0074 & Route acceleration & 4,5 / 3,4,5 & Reversal at $D=3$ s \\
\bottomrule
\end{tabularx}
\end{table}

표~\ref{tab:candidate-summary}에서 \texttt{scene-0028}의 반응은 첫 감속 판단 약 3.24초 뒤, 반복 판단 1.24초 뒤에 나타났다. \texttt{scene-0074}의 반응은 첫 가속 판단 약 3.36초 뒤, 반복 판단 0.86초 뒤에 나타났다. 제한 시간이 3초일 때 첫 판단 시각을 유지하면 두 장면 모두 실패하지만, 반복 판단에서 시계를 다시 시작하면 통과한다. 따라서 반복 사건에서 어느 시각부터 시간을 잴지 미리 정해야 한다. \texttt{scene-0042}와 \texttt{scene-0065}도 제한 시간 3초에서는 실패하고 4초부터 통과하여 제한 시간 선택의 영향을 보여 준다.

\begin{table}[H]
\centering
\bitablecaption{130개 사건에서 27개 속도 변화 판정 설정이 반응을 찾은 횟수.}{Number of response settings that find a speed change for each of 130 events.}
\label{tab:detector-sensitivity}
\footnotesize
\begin{tabular}{lr}
\toprule
Settings finding a response & Event count \\
\midrule
27/27 & 72 \\
18--26/27 & 51 \\
9--17/27 & 2 \\
1--8/27 & 4 \\
0/27 & 1 \\
\bottomrule
\end{tabular}
\end{table}

표~\ref{tab:detector-sensitivity}에서 72건은 27개 설정 모두가 반응을 찾았지만, 58건은 설정에 따라 판정이 달랐다. 한 후보는 어떤 설정에서도 반응을 찾지 못했고, 네 후보는 1--8개 설정에서만 반응을 찾았다. 즉 자차 속도에서 반응을 찾은 결과는 살펴보는 시간 구간, 최소 속도 변화와 요구 방향 일치 기준에 따라 달라진다. 따라서 검토 결과에는 사용한 탐지 설정도 함께 기록해야 한다.

\begin{table}[H]
\centering
\bitablecaption{연결 근거 없이 이전 장면의 요구를 넘기는 오류를 막는 시험.}{Test for blocking prior-requirement carryover without scene-continuity evidence.}
\label{tab:boundary-negative-control}
\scriptsize
\setlength{\tabcolsep}{2.5pt}
\begin{tabular}{lccccc}
\toprule
Case & Finished & Reset & No cross-scene & No carryover & Run \\
\midrule
correct & P & P & P & P & P \\
faulty transfer & P & F & F & F & P \\
\bottomrule
\end{tabular}
\end{table}

마지막으로 표~\ref{tab:boundary-negative-control}에서도 P와 F는 각각 통과와 실패를 뜻한다. 올바른 모델은 두 장면이 이어진다는 기록이 없는 경계에서 이전 요구를 폐기하고 모든 검사 조건을 통과했다. 오류를 넣은 모델은 첫 장면의 요구를 둘째 장면의 반응으로 잘못 완료했기 때문에, 경계에서 폐기해야 한다는 조건과 장면 사이 완료ㆍ이월을 금지한 조건을 위반했다. 따라서 자료 연결 근거가 없는 장면 결합을 막는 규칙이 의도대로 작동함을 확인했다.
